对称的运动方程并不保证最低能量态本身也具有同样的对称性。最简单的情形是势能最低点不是一个点，而是一整圈等能真空：理论允许沿这条圆周变换而不改变能量，但一旦实际选定其中一个真空，围绕该真空的“径向”与“沿圆周”涨落就不再等价。沿圆周移动没有恢复力，因此会出现一个无能隙的低能模。连续整体对称性自发破缺时出现这种无隙模的结论通常称为 Goldstone 定理；只从真空几何和守恒荷把这一结论推出来。

这里还必须区分两种性质不同的变换。若同一个变换参数在整个时空中取同一常数，它连接不同的物理态，是整体对称性；若参数可以逐点任意改变而只是在描述同一物理配置，它属于规范冗余。前一种情形的无能隙模是真实激发，后一种情形则会在规范场存在时与纵向规范自由度重新组合。整体 \(U(1)\) 复标量场已经足以显示“简并真空为何产生无能隙方向”；把同一对称性局域化以后，角向涨落才与纵向规范自由度重新组合。

采用的 SI 场归一化
\begin{equation*}
 [\varphi_{\rm ssb}]=\sqrt{\mathrm{J/m}},
\end{equation*}
使 \((\partial_\mu\varphi_{\rm ssb})^*(\partial^\mu\varphi_{\rm ssb})\) 具有 \(\mathrm{J/m^3}\) 量纲。取
\begin{equation}
 \Lag
 =(\partial_\mu\varphi_{\rm ssb})^*(\partial^\mu\varphi_{\rm ssb})-V(\varphi_{\rm ssb}),
 \qquad
 V(\varphi_{\rm ssb})
 =-\mu_{\rm ssb}^2\varphi_{\rm ssb}^*\varphi_{\rm ssb}
 +\frac{\lambda}{\hbar c}(\varphi_{\rm ssb}^*\varphi_{\rm ssb})^2,
 \label{eq:ssb-complex-scalar-lag-dp11}
\end{equation}
其中 \(\mu_{\rm ssb}\,[\mathrm{m^{-1}}]\)、\(\lambda>0\) 无量纲。势能只依赖 \(\varphi_{\rm ssb}^*\varphi_{\rm ssb}\)，故在 \(\varphi_{\rm ssb}\to \ee^{\ii\alpha}\varphi_{\rm ssb}\) 下不变。

非平凡真空由势能极小条件确定。定义实真空尺度 \(v\) 使
\begin{equation}
 \boxed{
 \langle\varphi_{\rm ssb}\rangle=\frac{v}{\sqrt2},
 \qquad
 v^2=\frac{\mu_{\rm ssb}^2\hbar c}{\lambda}.}
 \label{eq:ssb-vev-general-dp11}
\end{equation}
所有 \((v/\sqrt2)\ee^{\ii\theta}\) 具有相同能量；选择其中一个真空破坏整体 \(U(1)\)，但作用量仍保持该对称性。

在所选真空附近写成 Cartesian 涨落
\begin{equation}
 \varphi_{\rm ssb}(x)=\frac1{\sqrt2}\left[v+h(x)+\ii\pi(x)\right].
 \label{eq:ssb-cartesian-fields-dp11}
\end{equation}
势能 Hessian 的二次项为
\begin{equation}
 V^{(2)}=\mu_{\rm ssb}^2h^2+0\times\pi^2,
 \label{eq:ssb-quadratic-potential-dp11}
\end{equation}
与实标量质量项 \(\tfrac12(m^2c^2/\hbar^2)\phi^2\) 比较可得
\begin{equation}
 \boxed{
 m_h^2=\frac{2\hbar^2\mu_{\rm ssb}^2}{c^2},
 \qquad
 m_\pi=0.}
 \label{eq:ssb-higgs-goldstone-masses-dp11}
\end{equation}
\(h\) 沿真空流形的径向振动，\(\pi\) 沿连续简并真空的切向振动。质量为零不是参数偶合，而是势能沿对称轨道恒定所强制的 Hessian 零本征值。

上面的 Hessian 计算依赖具体的四次势，但“沿被破缺对称方向没有能隙”并不依赖这个势的特殊形式。对任意连续整体对称性，Noether 流 \(j_a^\mu\) 定义守恒荷；若这个荷作用在某个局域场上能把真空推向一个不同真空，就有
\begin{equation}
\boxed{
 Q_a=\int\dd^3 x\,j_a^0(x),
 \qquad
 \frac{\dd Q_a}{\dd t}=0,
 \qquad
 \langle\Omega|[Q_a,\phi_i(0)]|\Omega\rangle\neq0.}
\label{eq:goldstone-broken-charge-dp68}
\end{equation}
最后一个非零对易子就是“真空没有保持该生成元”的可计算判据。把完整能量--动量本征态插入这个对易子后，空间积分只留下零三动量中间态；而守恒荷又要求整个矩阵元不随时间振荡。因此谱中必须存在可以把能量连续压到零的激发，记为
\begin{equation}
\boxed{
 \exists\,|G_a(\mathbf p)\rangle:
 \qquad
 \lim_{\mathbf p\to0}E_{G_a}(\mathbf p)=0.}
\label{eq:goldstone-gapless-spectrum-dp68}
\end{equation}
Goldstone 定理所需的关键输入因此是作用在物理 Hilbert 空间上的整体守恒荷。局域规范变换属于描述冗余，不满足这一前提，必须另行分析其约束与物理自由度。

局域规范对称性的质量结构直接由一般群表示和标量真空的稳定子确定。设紧致规范群的一个简单因子为 $G$，标量多重态 $\varphi_{\rm ssb}$ 属于酉表示 $R$，其厄米生成元为 $T_R^a$，并定义
\begin{equation}
 D_\mu\varphi_{\rm ssb}
 =\left(\partial_\mu+\ii g_{\rm G}\mathscr C_\mu^aT_R^a\right)\varphi_{\rm ssb}.
 \label{eq:general-ssb-covariant-derivative-dp73}
\end{equation}
若势能的平稳真空为 $\langle\varphi_{\rm ssb}\rangle=v_R$，则在真空上
\begin{equation*}
 D_\mu v_R=\ii g_{\rm G}\mathscr C_\mu^aT_R^av_R.
\end{equation*}
因此标量动能的规范场二次项为
\begin{align*}
 (D_\mu v_R)^\dagger(D^\mu v_R)
 &=g_{\rm G}^2\mathscr C_\mu^a\mathscr C^{b\mu}
 v_R^\dagger T_R^aT_R^bv_R,\\
 &=\frac{g_{\rm G}^2}{2}\mathscr C_\mu^a\mathscr C^{b\mu}
 v_R^\dagger\{T_R^a,T_R^b\}v_R,
\end{align*}
其中第二行只用了 $\mathscr C_\mu^a\mathscr C^{b\mu}$ 对 $a,b$ 的对称性。与几何规范场归一化下的质量项
\[
\frac{\hbar c}{2}\mathscr C_\mu^a(M_g^2)_{ab}\mathscr C^{b\mu}
\]
比较，得到一般规范玻色子几何质量平方矩阵
\begin{equation}
 \boxed{
 (M_g^2)_{ab}
 =\frac{g_{\rm G}^2}{\hbar c}
 v_R^\dagger\{T_R^a,T_R^b\}v_R.}
 \label{eq:general-gauge-boson-mass-matrix-dp11}
\end{equation}
其本征值满足 $\lambda_r(M_g^2)=(m_rc/\hbar)^2$。更直接地，对任意实伴随向量 $x_a$，
\begin{equation}
 \boxed{
 x_a(M_g^2)_{ab}x_b
 =\frac{2g_{\rm G}^2}{\hbar c}
 \bigl\|x_aT_R^av_R\bigr\|^2\ge0.}
 \label{eq:gauge-mass-positive-quadratic-form-dp68}
\end{equation}
因此 $M_g^2$ 正半定，而零空间恰好由满足 $x_aT_R^av_R=0$ 的 Lie 代数组合组成；这些方向保持真空不变，构成未破缺子代数。非零本征方向对应被真空选择破缺的生成元，并产生有质量规范玻色子。这个结论完全由一般表示和真空决定。

现在才取单生成元的 $U(1)$ 特例，以显式看清自由度如何重组。令标量的所选 $U(1)$ 生成元本征值吸收到 $g_{\rm G}$ 中，采用几何规范联络 $\mathscr C_\mu\,[\mathrm{m^{-1}}]$，
\begin{equation*}
 D_\mu=\partial_\mu+\ii g_{\rm G}\mathscr C_\mu,
 \qquad
 \mathcal F_{\mu\nu}
 =\partial_\mu\mathscr C_\nu-\partial_\nu\mathscr C_\mu.
\end{equation*}
带有前述自发破缺复标量的这个阿贝尔规范模型通常称为 Abelian Higgs 模型，其拉格朗日量为
\begin{equation}
 \Lag_{\rm AH}
 =(D_\mu\varphi_{\rm ssb})^*(D^\mu\varphi_{\rm ssb})-V(\varphi_{\rm ssb})
 -\frac{\hbar c}{4}\mathcal F_{\mu\nu}\mathcal F^{\mu\nu}.
 \label{eq:abelian-higgs-lag-dp11}
\end{equation}
用极坐标参数化
\begin{equation}
 \varphi_{\rm ssb}(x)=\frac{v+h(x)}{\sqrt2}
 \exp\!\left[\frac{\ii\pi(x)}{v}\right].
 \label{eq:abelian-higgs-polar-dp11}
\end{equation}
局域变换 $\varphi_{\rm ssb}\to \ee^{\ii\alpha(x)}\varphi_{\rm ssb}$ 对应
\begin{equation*}
 \pi\to\pi+v\alpha,
 \qquad
 \mathscr C_\mu\to\mathscr C_\mu-\frac1{g_{\rm G}}\partial_\mu\alpha,
\end{equation*}
所以组合
\begin{equation}
 \boxed{
 \mathcal B_\mu
 \equiv\mathscr C_\mu+\frac1{g_{\rm G}v}\partial_\mu\pi}
 \label{eq:gauge-invariant-massive-vector-dp11}
\end{equation}
在局域变换下不变。标量动能变为
\begin{equation}
 (D_\mu\varphi_{\rm ssb})^*(D^\mu\varphi_{\rm ssb})
 =\frac12(\partial h)^2
 +\frac{g_{\rm G}^2}{2}(v+h)^2\mathcal B_\mu\mathcal B^\mu.
 \label{eq:abelian-higgs-kinetic-expanded-dp11}
\end{equation}
其中二次项正是一般质量矩阵~\eqnrefs{eq:general-gauge-boson-mass-matrix-dp11} 在一维表示、$v_R=v/\sqrt2$ 时的特例。与
\[
 \frac{\hbar c}{2}
 \left(\frac{m_Ac}{\hbar}\right)^2
 \mathcal B_\mu\mathcal B^\mu
\]
比较得到
\begin{equation}
 \boxed{m_Ac^2=g_{\rm G}\sqrt{\hbar c}\,v.}
 \label{eq:abelian-higgs-vector-mass-dp11}
\end{equation}
若定义能量形式真空尺度 $v_E\equiv\sqrt{\hbar c}\,v$，则 $m_Ac^2=g_{\rm G}v_E$。破缺前的两个实标量自由度与两个无质量规范横向偏振总数为 4；破缺后重组为一个径向标量和三个有质量矢量偏振，自由度总数不变。

微扰计算若保留 $\mathscr C_\mu$ 与 $\pi$ 为独立场，需要选择能消去二次混合的规范固定。Cartesian 参数化产生
\begin{equation}
 \Lag^{(2)}\supset g_{\rm G}v\,\mathscr C^\mu\partial_\mu\pi.
 \label{eq:gauge-goldstone-mixing-dp11}
\end{equation}
定义
\begin{equation}
 \boxed{
 F=\partial^\mu\mathscr C_\mu
 -\xi\frac{g_{\rm G}v}{\hbar c}\,\pi,}
 \label{eq:rxi-gauge-functional-abelian-dp42}
\end{equation}
并加入
\begin{equation}
 \Lag_{\rm gf}
 =-\frac{\hbar c}{2\xi}F^2.
 \label{eq:rxi-abelian-gf-dp11}
\end{equation}
这种一参数规范选择通常称为 $R_\xi$ 规范。其交叉项经分部积分后抵消\eqnrefs{eq:gauge-goldstone-mixing-dp11}，剩余 $\pi^2$ 项给出
\begin{equation}
 \boxed{m_\pi^2=\xi m_A^2.}
 \label{eq:rxi-goldstone-mass-abelian-dp42}
\end{equation}
因此 $\pi$ 在固定规范中的质量依赖 $\xi$，不能把它解释成独立可观测粒子的质量；可观测散射矩阵中的规范参数依赖必须在完整规范恒等式中相消。

对外部有质量规范玻色子的物理四动量 \(k^\mu\)，若 \(E\gg m_Vc^2\)，纵向偏振满足
\begin{equation*}
 \epsilon_L^\mu(k)
 =\frac{k^\mu}{m_Vc}
 +\order\!\left(\frac{m_Vc^2}{E}\right).
\end{equation*}
Slavnov--Taylor 恒等式把 \(k_\mu\mathcal M^\mu\) 与对应 Goldstone 振幅联系起来，因此
\begin{equation}
 \mathcal M(\cdots V_L\cdots)
 =C\,\mathcal M(\cdots\pi\cdots)
 +\order\!\left(\frac{m_Vc^2}{E}\right),
 \label{eq:goldstone-equivalence-general-dp11}
\end{equation}
其中常用壳上树级约定给 \(C=1\)。控制量为 \(\epsilon_{\rm GE}=m_Vc^2/E\)。当 \(\epsilon_{\rm GE}\not\ll1\)（尤其接近产生阈值）时，高能展开失效，必须使用完整纵向偏振矢量与规范一致的原振幅；高阶修正还会改变有限的匹配因子 \(C\)。


整体破缺生成元通过守恒荷要求真实无质量 Goldstone 模；局域规范冗余经规范固定与 BRST 量子化后，把真空切向自由度与纵向规范偏振重组为有质量矢量表示。一般质量矩阵\eqnrefs{eq:general-gauge-boson-mass-matrix-dp11} 的零空间给出未破缺子代数，非零本征值给出规范玻色子质量。






对于一般多分量标量场，前面的圆形真空可以写成一个不依赖具体表示的代数结论。若势在生成元 \(T^a\) 下不变，则
\begin{equation}
\boxed{
\frac{\partial V}{\partial\phi_i}(T^a)_{ij}\phi_j=0.}
\label{eq:ssb-general-potential-invariance-dp68}
\end{equation}
在平稳真空 \(v\) 处再对场求导，定义 Hessian \(H_{ki}(v)=\partial_k\partial_iV|_v\)，得到
\begin{equation}
\boxed{
H_{ki}(v)(T^av)_i=0.}
\label{eq:ssb-general-hessian-zero-mode-dp68}
\end{equation}
因此只有当 \(T^av\neq0\) 时，这个生成元才给出一个非平凡 Hessian 零方向。


\begin{derivation}[从\eqnrefs{eq:ssb-general-potential-invariance-dp68}到\eqnrefs{eq:ssb-general-hessian-zero-mode-dp68}]
把一组实标量场坐标记为 $\phi_i$。连续整体对称性的无穷小作用写成

\begin{equation*}
 \delta\phi_i=\epsilon^a(T^a)_{ij}\phi_j,
\end{equation*}
其中 $\epsilon^a$ 是无穷小实参数，$T^a$ 是该实场表示中的生成元。势能在整个场空间上保持不变，因此对任意场构型都有
\begin{align*}
 0=\delta V
 &=\frac{\partial V}{\partial\phi_i}\,\delta\phi_i,\\
 &=\epsilon^a\frac{\partial V}{\partial\phi_i}(T^a)_{ij}\phi_j.
\end{align*}
由于各 $\epsilon^a$ 相互独立，每个生成元分别满足函数恒等式
\begin{equation*}
 \frac{\partial V}{\partial\phi_i}(T^a)_{ij}\phi_j=0.
\end{equation*}
对 $\phi_k$ 再求一次偏导，乘积法则给
\begin{equation*}
 0=
 \frac{\partial^2V}{\partial\phi_k\partial\phi_i}
 (T^a)_{ij}\phi_j
 +\frac{\partial V}{\partial\phi_i}(T^a)_{ik}.
\end{equation*}
令 $v_i$ 为平稳真空，即
\begin{equation*}
 \left.\frac{\partial V}{\partial\phi_i}\right|_{\phi=v}=0.
\end{equation*}
第二项在真空处消失，于是
\begin{align*}
 0
 &=\left.\frac{\partial^2V}{\partial\phi_k\partial\phi_i}
 \right|_{v}(T^av)_i,\\
 &=H_{ki}(v)(T^av)_i,
\end{align*}
其中
\begin{equation*}
 H_{ki}(v)\equiv
 \left.\frac{\partial^2V}{\partial\phi_k\partial\phi_i}\right|_v
\end{equation*}
是势能在真空处的 Hessian 矩阵。若生成元 $T^a$ 被真空自发破缺，则 $T^av\neq0$；上式说明非零向量 $T^av$ 属于 $H(v)$ 的零空间。对具有标准正定动能项的实标量场，二次拉氏量中的质量平方矩阵正比于该 Hessian 矩阵，所以这个真空轨道切向方向对应零质量模。Goldstone 零模因此来自“势能沿连续对称轨道保持常数”这一一般几何事实，而不依赖墨西哥帽势的具体形状。
\end{derivation}

\begin{derivation}[从\eqnrefs{eq:general-gauge-boson-mass-matrix-dp11}到\eqnrefs{eq:gauge-mass-positive-quadratic-form-dp68}]

正文\eqnrefs{eq:general-gauge-boson-mass-matrix-dp11} 给出生成元空间中的质量矩阵。要判断它是否可能产生负质量平方，不需要先对角化；只需把它与任意实向量 \(x_a\) 收缩，并把结果改写成真空向量的范数平方。


定义
\begin{equation*}
 (M_g^2)_{ab}
 =\frac{g_{\rm G}^2}{\hbar c}
 v_R^\dagger\{T^a,T^b\}v_R.
\end{equation*}
对任意实伴随空间向量 $x_a$，令
\begin{equation*}
 X\equiv x_aT^a.
\end{equation*}
若采用厄米生成元 $T^{a\dagger}=T^a$，则 $X^\dagger=X$。质量矩阵对应的二次型逐步化为
\begin{align*}
 x_a(M_g^2)_{ab}x_b
 &=\frac{g_{\rm G}^2}{\hbar c}
 v_R^\dagger x_ax_b\{T^a,T^b\}v_R,\\
 &=\frac{g_{\rm G}^2}{\hbar c}
 v_R^\dagger\{X,X\}v_R,\\
 &=\frac{2g_{\rm G}^2}{\hbar c}v_R^\dagger X^2v_R,\\
 &=\frac{2g_{\rm G}^2}{\hbar c}(Xv_R)^\dagger(Xv_R),\\
 &=\frac{2g_{\rm G}^2}{\hbar c}\|Xv_R\|^2\ge0.
\end{align*}
因此 $M_g^2$ 是实对称正半定矩阵，树级谱不会从这一机制产生负质量平方的矢量模。

等号成立当且仅当
\begin{equation*}
 Xv_R=0.
\end{equation*}
于是质量矩阵的零空间恰由所有湮灭真空的 Lie 代数方向组成，也就是未破缺子代数。反之，若 $Xv_R\neq0$，对应二次型严格为正，该生成元方向属于破缺子空间，并为相应规范玻色子提供非零质量。破缺方向上的 Goldstone 切向自由度与规范场纵向偏振重组为有质量矢量表示；这一线性代数结论不依赖某个具体规范群的显式矩阵对角化。
\end{derivation}


\begin{derivation}[从\eqnrefs{eq:ssb-complex-scalar-lag-dp11}到\eqnrefs{eq:ssb-higgs-goldstone-masses-dp11}]

正文\eqnrefs{eq:ssb-complex-scalar-lag-dp11} 给出了只依赖 \(|\varphi_{\rm ssb}|\) 的对称势。先由驻点条件确定真空半径，再在所选真空附近展开到二次阶；径向与切向二次系数的差别最终给出正文\eqnrefs{eq:ssb-higgs-goldstone-masses-dp11} 的有质量径向模和无质量切向模。


令 \(r^2\equiv\varphi_{\rm ssb}^*\varphi_{\rm ssb}\)，则
\begin{equation*}
 V(r)=-\mu_{\rm ssb}^2r^2+\frac\lambda{\hbar c}r^4.
\end{equation*}
极值条件为
\begin{align*}
 \frac{\dd V}{\dd r}
 &=-2\mu_{\rm ssb}^2r+\frac{4\lambda}{\hbar c}r^3,\\
 &=2r\left(-\mu_{\rm ssb}^2+\frac{2\lambda}{\hbar c}r^2\right)=0.
\end{align*}
除去不稳定的 \(r=0\)，稳定极小值满足
\begin{equation*}
 r^2=\frac{\mu_{\rm ssb}^2\hbar c}{2\lambda}.
\end{equation*}
定义 \(r=v/\sqrt2\) 即得到
\begin{equation*}
 v^2=\frac{\mu_{\rm ssb}^2\hbar c}{\lambda}.
\end{equation*}

将真空附近参数化代入
\begin{equation*}
 \varphi_{\rm ssb}=\frac1{\sqrt2}(v+h+\ii\pi),
 \qquad
 \varphi_{\rm ssb}^*\varphi_{\rm ssb}=\frac12[(v+h)^2+\pi^2].
\end{equation*}
保留到二次阶：
\begin{align*}
 V
 ={}&-\frac{\mu_{\rm ssb}^2}{2}
 (v^2+2vh+h^2+\pi^2)\\
 &+\frac\lambda{4\hbar c}
 (v^2+2vh+h^2+\pi^2)^2.
\end{align*}
线性 \(h\) 项系数为
\begin{equation*}
 -\mu_{\rm ssb}^2v+\frac{\lambda v^3}{\hbar c}=0,
\end{equation*}
二次项为
\begin{equation*}
 V^{(2)}=\mu_{\rm ssb}^2h^2+0\times\pi^2.
\end{equation*}
与 \(\tfrac12(m^2c^2/\hbar^2)\phi^2\) 比较即得\eqnrefs{eq:ssb-higgs-goldstone-masses-dp11}。
\end{derivation}

\begin{derivation}[从\eqnrefs{eq:goldstone-broken-charge-dp68}到\eqnrefs{eq:goldstone-gapless-spectrum-dp68}]
设连续整体对称性有守恒流 \(j_a^\mu\) 与荷

\begin{equation*}
 Q_a=\int\dd^3 x\,j_a^0(x),
 \qquad
 \frac{\dd Q_a}{\dd t}=0.
\end{equation*}
若某个局域场满足
\begin{equation*}
 \langle\Omega|[Q_a,\phi_i(0)]|\Omega\rangle\neq0,
\end{equation*}
则把 \(Q_a\) 写成空间积分并在对易子中插入完整能量--动量本征态，得到
\begin{align*}
 \langle[Q_a,\phi_i(0)]\rangle
 =\int\dd^3 x\sum_n
 \bigl[&\langle\Omega|j_a^0(x)|n\rangle
 \langle n|\phi_i(0)|\Omega\rangle\\
 &-\langle\Omega|\phi_i(0)|n\rangle
 \langle n|j_a^0(x)|\Omega\rangle\bigr].
\end{align*}
平移不变性给
\begin{equation*}
 \langle\Omega|j_a^0(x)|n\rangle
 =\ee^{-\ii p_n\cdot x/\hbar}
 \langle\Omega|j_a^0(0)|n\rangle.
\end{equation*}
空间积分产生 \((2\pi\hbar)^3\delta^{(3)}(\mathbf p_n)\)，故只有零三动量态贡献。左边又因 \(Q_a\) 守恒与真空定态性而与时间无关；若所有零三动量中间态都有严格正的能隙，则右边只能含非零频率相位，不能维持非零常数。因此谱必须在 \(\mathbf p\to0\) 时具有 \(E\to0\) 的支撑，即 Goldstone 模。
\end{derivation}

\begin{derivation}[从\eqnrefs{eq:abelian-higgs-polar-dp11}到\eqnrefs{eq:abelian-higgs-vector-mass-dp11}]
从
\[
\varphi_{\rm ssb}(x)
=\frac{v+h(x)}{\sqrt2}\ee^{\ii\pi(x)/v},
\qquad
D_\mu=\partial_\mu+\ii g_{\rm G}\mathscr C_\mu
\]
出发。局域 $U(1)$ 变换
\[
\varphi_{\rm ssb}\to \ee^{\ii\alpha(x)}\varphi_{\rm ssb},
\qquad
\mathscr C_\mu\to
\mathscr C_\mu-\frac1{g_{\rm G}}\partial_\mu\alpha
\]
要求极坐标变量变成
\[
h\to h,
\qquad
\pi\to\pi+v\alpha.
\]
因此组合
\[
\mathcal B_\mu
\equiv
\mathscr C_\mu+\frac1{g_{\rm G}v}\partial_\mu\pi
\]
满足 $\mathcal B_\mu\to\mathcal B_\mu$，它才是破缺相位中的规范不变矢量变量。

显式求协变导数：
\begin{align*}
D_\mu\varphi_{\rm ssb}
&=\frac{\ee^{\ii\pi/v}}{\sqrt2}
\left[
\partial_\mu h
+\ii(v+h)
\left(\frac{\partial_\mu\pi}{v}+g_{\rm G}\mathscr C_\mu\right)
\right]\\
&=\frac{\ee^{\ii\pi/v}}{\sqrt2}
\left[
\partial_\mu h
+\ii g_{\rm G}(v+h)\mathcal B_\mu
\right].
\end{align*}
取模平方时，实的径向项与纯虚的相位项交叉项严格相消：
\[
(D_\mu\varphi)^*(D^\mu\varphi)
=\frac12(\partial h)^2
+\frac{g_{\rm G}^2}{2}(v+h)^2\mathcal B_\mu\mathcal B^\mu.
\]
展开 $(v+h)^2=v^2+2vh+h^2$。与矢量二次型有关的纯二次项是
\[
\frac{g_{\rm G}^2v^2}{2}\mathcal B_\mu\mathcal B^\mu.
\]
在上述 SI 归一化下，Proca 质量项写成
\[
\frac{\hbar c}{2}
\left(\frac{m_Ac}{\hbar}\right)^2
\mathcal B_\mu\mathcal B^\mu.
\]
两者是同一个二次算符，所以系数必须相等：
\[
g_{\rm G}^2v^2
=\hbar c\left(\frac{m_Ac}{\hbar}\right)^2,
\]
从而
\[
\boxed{m_Ac^2=g_{\rm G}\sqrt{\hbar c}\,v.}
\]
Goldstone 场 $\pi$ 通过 $\mathcal B_\mu$ 重组为有质量矢量场的纵向偏振。破缺前自由度计数为“复标量的两个实自由度 + 无质量矢量的两个横向偏振”共四个；破缺后为“径向标量 $h$ 一个 + 有质量矢量 $\mathcal B_\mu$ 三个偏振”仍共四个。这个计数是上述变量重组没有丢失物理自由度的独立检查。
\end{derivation}

\begin{derivation}[从\eqnrefs{eq:gauge-goldstone-mixing-dp11}到\eqnrefs{eq:rxi-goldstone-mass-abelian-dp42}]
Cartesian 展开给出
\[
\Lag_{\rm mix}=g_{\rm G}v\,\mathscr C^\mu\partial_\mu\pi,
\qquad
F=\partial\!\cdot\!\mathscr C-\xi\frac{g_{\rm G}v}{\hbar c}\pi .
\]
规范固定项展开为
\[
\Lag_{\rm gf}
=-\frac{\hbar c}{2\xi}(\partial\!\cdot\!\mathscr C)^2
+g_{\rm G}v\,\pi\,\partial\!\cdot\!\mathscr C
-\frac{\xi g_{\rm G}^2v^2}{2\hbar c}\pi^2 .
\]
对中间项分部积分，
\[
g_{\rm G}v\,\pi\,\partial_\mu\mathscr C^\mu
=-g_{\rm G}v\,\mathscr C^\mu\partial_\mu\pi+\text{边界项},
\]
故它与 $\Lag_{\rm mix}$ 相消。剩余 $\pi^2$ 项与标准质量项比较：
\[
-\frac{\xi g_{\rm G}^2v^2}{2\hbar c}\pi^2
=-\frac12\left(\frac{m_\pi c}{\hbar}\right)^2\pi^2,
\qquad
\left(\frac{m_Ac}{\hbar}\right)^2=\frac{g_{\rm G}^2v^2}{\hbar c},
\]
于是
\[
\boxed{m_\pi^2=\xi m_A^2}.
\]
$\pi$ 的质量依赖规范参数 $\xi$，而物理规范玻色子质量 $m_A$ 不随 $\xi$ 改变。
\end{derivation}

